\chapter{Antecedentes}\label{Chap:Antecedentes}

Esta línea de investigación continúa dos trabajos previos que abordaron, en dispositivos distintos, el problema que se estudia aquí: el control de posición en presencia de zona muerta y fricción.

\section{Control del ECP-730 en configuración repulsiva}

Hernández Alcántara \cite{diana_tesis} estudió el sistema de levitación magnética ECP-730 en su configuración repulsiva, en la que la bobina inferior empuja al imán hacia arriba y el sistema es estable en lazo abierto. El trabajo obtuvo un modelo no lineal del dispositivo, formado por un sistema lineal de segundo orden y una función no lineal que multiplica la entrada, y lo validó con experimentos en tiempo real. A partir de ensayos de estado estacionario y de respuesta dinámica identificó los parámetros $a$, $b$ y $c_1$ que se emplean en esta tesis. Debido a la dificultad de modelar con precisión la fricción, el modelo solo incluyó su componente viscosa; las diferencias entre el modelo y el dispositivo real se atribuyeron a fenómenos no modelados, en particular la fricción y la zona muerta. De ahí se concluyó que, para reducir sus efectos, el sistema de control debe proporcionar alta ganancia en estado estacionario.

El controlador propuesto combina una linealización por retroalimentación de estados, que cancela la no linealidad de la fuerza electromagnética y ubica los polos del lazo cerrado, con un lazo externo de estructura variable. Como no todos los estados se miden directamente, se diseñaron e implementaron observadores lineal y no lineal. Las pruebas en tiempo real mostraron la efectividad del esquema en regulación y seguimiento. En un trabajo posterior \cite{alcantara}, el lazo externo se diseñó con acción integral en el dominio de la frecuencia. Con él se obtuvieron mejores respuestas transitorias y menores errores en estado estacionario, pero la interacción entre la acción integral y la fricción entre el imán, en forma de disco, y la varilla de vidrio que lo guía provocó atascamiento-deslizamiento, con oscilaciones de tipo ciclo límite. Estas oscilaciones se eliminaron con una estrategia de control conmutado basada en la caracterización experimental de la superficie de deslizamiento, que conserva un error bajo en estado estacionario.

\section{Control con doble acción integral de un motor de CD}

Pérez-Gómez \cite{perez} estudió el control de posición de un motor de corriente directa (CD) de imanes permanentes con zona muerta y fricción. El trabajo desarrolló un modelo no lineal del motor, con una representación de la zona muerta llamada ``dura'' ubicada entre los subsistemas eléctrico y mecánico, y lo validó con el motor real a bajas velocidades. Sobre este modelo comparó tres estrategias: un controlador PI con cancelación mediante zona muerta inversa, un controlador conmutado PI--PD y un controlador lineal con doble acción integral y compensador de adelanto (PII). La cancelación con zona muerta inversa produjo oscilaciones de alta frecuencia (\emph{chattering}), y el control conmutado ofreció un error en estado estacionario muy bajo y sin oscilaciones, aunque su análisis, diseño e implementación resultan complejos. El PII, en cambio, se implementó en tiempo real y mostró un error en estado estacionario y un sobreimpulso bajos, con la ventaja de poder analizarse con las herramientas del control clásico y de ser sencillo de diseñar e implementar.

\section{Relación con este trabajo}

Ambos trabajos validaron sus controladores en dispositivos estables en lazo abierto. Esta tesis toma el modelo y los parámetros del primero y la estructura de control del segundo, y los aplica a un caso que ninguno abordó: el ECP-730 en configuración atractiva, inestable en lazo abierto, en el que el controlador debe estabilizar el sistema y, al mismo tiempo, lidiar con el atascamiento-deslizamiento.
